\section{Check a format whose fields agree}
The format payload identifies sixteen-bit integer PCM with format tag one. The worked
values are channels two, frame rate 16000, byte rate 64000, frame size four and sample
width sixteen bits. A sixteen-byte format payload holds these six fields. The local
reader also accepts an eighteen-byte payload whose final two extension-size bytes are
zero; it does not interpret any other extension.

Offsets in this table begin at the format payload, which is file offset 32 in the
worked input. Each numerical field is unsigned and little endian.
\tagpdfsetup{table/header-rows={1}}
\begin{center}
\begin{tabular}{rlrr}
Payload offset & Field & Bytes & Worked value \\
\hline
0 & Format tag & 2 & 1 \\
2 & Channel count & 2 & 2 \\
4 & Frame rate & 4 & 16000 \\
8 & Byte rate & 4 & 64000 \\
12 & Frame size & 2 & 4 \\
14 & Bits per sample & 2 & 16 \\
16 & Extension size, if present & 2 & 0 \\
\end{tabular}
\end{center}
The last row is absent from the worked sixteen-byte payload. Test its value only
when the validated payload length is eighteen; otherwise those two bytes belong
to the following chunk header.

Derive the redundant quantities after validating the independent ones. With $C$
channels and sixteen-bit codes, one frame has $2C$ bytes. At rate $F$, the required byte
rate is $2CF$. Compare the stored frame size and byte rate with those results. A
stereo declaration paired with frame size two is inconsistent even if the data length
is even. The local bounds keep the correct products small, but perform their formation
without accidentally retaining a signed field interpretation.

Keep the distinction between format payload and compiler structure layout. The file's
six fields have specified byte widths and positions; a native language structure may
include fields or alignment rules that the file does not contain. Assemble unsigned
fields from their actual bytes. The PCM sample codes use a signed interpretation,
which is a later operation on a different region.
